\section*{ID6585: Why Use \textbackslash{}(32/27\textbackslash{}) and How to Make It Fractal: R := (27)/(32) ≈ 0.84375.}
\paragraph{Metadata.}
PiID: 5821167456270 | RTEID: RTE:P25571.0174:T4.4341 | UUID: 85e56962-8200-5b20-875c-5309c1d6d35f

\paragraph{Canonical Statement.}
I prefer the first: set the linear contraction ratio \textbackslash{}[ r := \textbackslash{}frac\{27\}\{32\} \textbackslash{}approx 0.84375. \textbackslash{}] Use the kissing number \textbackslash{}(N\textbackslash{}) of the chosen geometry as the number of copies produced at each step. That gives a canonical self-similar IFS.

\paragraph{Canonical Equation.}
\[
r := \frac{27}{32} \approx 0.84375.
\]

\paragraph{Dictionary.}
I prefer the first: set the linear contraction ratio Use the kissing number \textbackslash{}(N\textbackslash{}) of the chosen geometry as the number of copies produced at each step.

\paragraph{Encyclopedia.}
Why Use \textbackslash{}(32/27\textbackslash{}) and How to Make It Fractal: R := (27)/(32) ≈ 0.84375. is a scaling / order-of-magnitude relation, written r := (27)/(32) ≈ 0.84375.. It expresses a scaling or proportionality, capturing how the quantity grows with its parameters. It is built from r. Within the Trawin Zero Point registry it serves as one element of the framework's mathematical narrative.
